\documentclass[11pt,a4paper]{article}
\usepackage[utf8]{inputenc}
\usepackage[french]{babel}
\usepackage{amsmath,amssymb,amsthm}
\usepackage{geometry}
\geometry{hmargin=2.5cm,vmargin=3cm}
\usepackage{tikz}
\usetikzlibrary{cd,positioning,shapes}
\usepackage{listings}
\usepackage{xcolor}

\lstset{
    language=Python,
    basicstyle=\ttfamily\small,
    keywordstyle=\color{blue}\bfseries,
    stringstyle=\color{red},
    commentstyle=\color{gray}\itshape,
    numbers=left,
    numberstyle=\tiny\color{gray},
    breaklines=true,
    frame=single,
    showstringspaces=false
}

\title{\Huge INTRODUCTION AUX DISTRIBUTIONS DE SCHWARTZ \\ \large Cours magistral, exercices corrigés et travaux pratiques}
\author{\Large Charles EDOU NZE}
\date{\today}

\begin{document}
\maketitle
\tableofcontents
\newpage

\part{Cours Magistral}

\section*{Jalon 82 : Introduction aux distributions de Schwartz}

\section{Introduction : Genèse physique et mathématique}

L'étude des phénomènes physiques impulsifs ou ponctuels a longtemps posé un défi majeur à l'analyse mathématique classique. Considérons, par exemple, le choc instantané d'un marteau sur une corde vibrante, ou encore la modélisation d'une charge électrique ponctuelle. Dans ces situations, on observe une quantité totale non nulle (comme une impulsion ou une charge) concentrée en un unique point de l'espace ou du temps.

Si l'on cherche à représenter une telle densité par une fonction classique $f(x)$, cette fonction devrait être nulle partout sauf en ce point précis (disons, l'origine $x=0$), où elle devrait tendre vers l'infini. Or, l'intégrale de Lebesgue d'une fonction nulle presque partout est nulle, ce qui contredit l'idée que l'intégrale totale représente la charge globale (qui vaut 1). Les ingénieurs et physiciens, tels que Paul Dirac au début du XXe siècle, utilisaient des objets formels comme la "fonction" $\delta$ (qui porte désormais son nom), la manipulant avec succès dans des calculs sans justification mathématique rigoureuse.

La théorie des distributions, développée par le mathématicien français Laurent Schwartz (médaille Fields 1950), apporte une solution élégante et définitive à ce paradoxe. L'idée fondatrice est de changer de perspective : au lieu de considérer une "fonction" ponctuellement, on l'étudie par son interaction (son intégrale) avec des fonctions dites "tests", très régulières. Les objets mathématiques ne sont plus étudiés par leurs valeurs ponctuelles, mais par la manière dont ils "mesurent" d'autres fonctions. Cette approche généralise la notion de fonction et donne un cadre rigoureux à la masse de Dirac et à ses dérivées.

\begin{center}
\begin{tikzpicture}[scale=1]
  \draw[->] (-3,0) -- (3,0) node[right] {$x$};
  \draw[->] (0,-0.5) -- (0,4) node[above] {$y$};

  % Suite de fonctions approchant Dirac
  \draw[blue, thick] (-2,0) -- (-0.5,0) -- (-0.5, 1) -- (0.5, 1) -- (0.5, 0) -- (2,0);
  \node[blue] at (1, 1.2) {$f_1$};

  \draw[red, thick] (-2,0) -- (-0.25,0) -- (-0.25, 2) -- (0.25, 2) -- (0.25, 0) -- (2,0);
  \node[red] at (0.6, 2.2) {$f_2$};

  \draw[green!60!black, thick] (-2,0) -- (-0.125,0) -- (-0.125, 4) -- (0.125, 4) -- (0.125, 0) -- (2,0);
  \node[green!60!black] at (0.4, 3.8) {$f_3$};

  \node[below] at (0,0) {$0$};
\end{tikzpicture}
\\ \textit{Figure 1 : Fonctions "portes" d'aire 1 dont la limite au sens des distributions est la distribution de Dirac $\delta_0$.}
\end{center}

\section{Définitions, Théorèmes \& Exemples}

Pour définir le concept de distribution, il faut préalablement définir l'espace des fonctions sur lequel ces distributions vont agir. Ces fonctions doivent être suffisamment régulières pour absorber les irrégularités des objets qu'on étudie.

\begin{quote}\textbf{Définition 1 (Support d'une fonction) :} Soit \end{quote}$f : \mathbb{R}^n \to \mathbb{C}$ une fonction continue. Le support de $f$, noté $\text{supp}(f)$, est l'adhérence de l'ensemble des points où $f$ est non nulle :
\begin{quote}\end{quote}$$\text{supp}(f) = \overline{\{ x \in \mathbb{R}^n \mid f(x) \neq 0 \}}$$
\begin{quote}Une fonction est dite à \textbf{support compact} si son support est un sous-ensemble compact (fermé et borné) de \end{quote}$\mathbb{R}^n$.

\textbf{Exemple concret de calcul de support :}
Considérons la fonction en une dimension :
$f(x) = 1$ si $x \in ]-1, 1[$, et $f(x) = 0$ sinon.
L'ensemble des points où $f$ est non nulle est l'ouvert $]-1, 1[$.
L'adhérence de cet ouvert est le segment $[-1, 1]$.
Ainsi, $\text{supp}(f) = [-1, 1]$, qui est un compact (fermé et borné).

\begin{quote}**Définition 2 (Espace des fonctions tests \end{quote}$\mathcal{D}(\mathbb{R}^n)$) :**
\begin{quote}L'espace \end{quote}$\mathcal{D}(\mathbb{R}^n)$ (souvent noté $\mathcal{C}_c^\infty(\mathbb{R}^n)$) est l'espace vectoriel des fonctions $\phi : \mathbb{R}^n \to \mathbb{C}$ qui sont :
\begin{quote}1. Infiniment dérivables (de classe \end{quote}$\mathcal{C}^\infty$).
\begin{quote}2. À support compact.\end{quote}

Il n'est pas immédiat qu'il existe des fonctions non nulles dans $\mathcal{D}(\mathbb{R})$. Les polynômes, par exemple, sont $\mathcal{C}^\infty$ mais n'ont jamais un support compact (sauf le polynôme nul).

\textbf{Exemple fondamental de fonction test (la fonction plateau) :}
Soit $\phi(x) = \exp\left(-\frac{1}{1-x^2}\right)$ pour $|x| < 1$, et $\phi(x) = 0$ pour $|x| \ge 1$.
Cette fonction est infiniment dérivable sur tout $\mathbb{R}$ (en particulier en $-1$ et en $1$, car toutes ses dérivées tendent vers $0$) et son support est $[-1, 1]$, qui est compact. Ainsi, $\phi \in \mathcal{D}(\mathbb{R})$.

Avant de définir rigoureusement les distributions en tant que formes linéaires continues, nous devons définir la notion de convergence sur l'espace $\mathcal{D}(\mathbb{R}^n)$.

\begin{quote}**Définition 3 (Convergence dans \end{quote}$\mathcal{D}$) :**
\begin{quote}Une suite \end{quote}$(\phi_j)_{j \in \mathbb{N}}$ de fonctions de $\mathcal{D}(\mathbb{R}^n)$ converge vers une fonction $\phi \in \mathcal{D}(\mathbb{R}^n)$ si :
\begin{quote}1. Il existe un compact fixe \end{quote}$K \subset \mathbb{R}^n$ tel que, pour tout $j \in \mathbb{N}$, $\text{supp}(\phi_j) \subset K$ et $\text{supp}(\phi) \subset K$.
\begin{quote}2. Pour tout multi-indice \end{quote}$\alpha$, la suite des dérivées partielles $\left( \partial^\alpha \phi_j \right)$ converge uniformément vers $\partial^\alpha \phi$ sur $K$.

\begin{quote}\textbf{Définition 4 (Distribution de Schwartz) :} Une \textbf{distribution} sur \end{quote}$\mathbb{R}^n$ est une forme linéaire $T$ sur l'espace vectoriel $\mathcal{D}(\mathbb{R}^n)$, c'est-à-dire une application $T : \mathcal{D}(\mathbb{R}^n) \to \mathbb{C}$ telle que :
\begin{quote}1. \textbf{Linéarité :} \end{quote}$\forall \phi, \psi \in \mathcal{D}$, $\forall \lambda, \mu \in \mathbb{C}$, $T(\lambda \phi + \mu \psi) = \lambda T(\phi) + \mu T(\psi)$.
\begin{quote}2. \textbf{Continuité séquentielle :} Pour toute suite \end{quote}$(\phi_j)$ convergeant vers $0$ dans $\mathcal{D}$, on a $\lim_{j \to +\infty} T(\phi_j) = 0$.
\begin{quote}L'ensemble de toutes les distributions sur \end{quote}$\mathbb{R}^n$ est noté $\mathcal{D}'(\mathbb{R}^n)$. L'action d'une distribution $T$ sur une fonction test $\phi$ s'écrit souvent sous forme de crochet de dualité : $\langle T, \phi \rangle$.

\textbf{Exemple 1 : La distribution de Dirac (Distribution singulière)}
Soit $a \in \mathbb{R}^n$. On définit l'application $\delta_a : \mathcal{D}(\mathbb{R}^n) \to \mathbb{C}$ par :
$$\langle \delta_a, \phi \rangle = \phi(a)$$
Montrons que c'est une distribution.
- Linéarité : $\langle \delta_a, \lambda \phi + \mu \psi \rangle = (\lambda \phi + \mu \psi)(a) = \lambda \phi(a) + \mu \psi(a) = \lambda \langle \delta_a, \phi \rangle + \mu \langle \delta_a, \psi \rangle$.
- Continuité : Si $\phi_j \to 0$ dans $\mathcal{D}$, cela implique la convergence uniforme (en particulier ponctuelle) vers $0$, donc $\phi_j(a) \to 0$, d'où $\langle \delta_a, \phi_j \rangle \to 0$.
Ainsi, $\delta_a \in \mathcal{D}'(\mathbb{R}^n)$.

\textbf{Exemple 2 : Distributions régulières (généralisation des fonctions)}
Soit $f \in L^1_{loc}(\mathbb{R}^n)$, c'est-à-dire une fonction localement intégrable (son intégrale de Lebesgue est finie sur tout compact de $\mathbb{R}^n$). On lui associe $T_f$ par :
$$\langle T_f, \phi \rangle = \int_{\mathbb{R}^n} f(x) \phi(x) dx$$
- Linéarité : par la linéarité de l'intégrale.
- Continuité : Soit $\phi_j \to 0$ dans $\mathcal{D}$. Toutes les $\phi_j$ ont leur support inclus dans un compact $K$. Alors,
  $$\left| \langle T_f, \phi_j \rangle \right| \le \int_K |f(x)| |\phi_j(x)| dx \le \|\phi_j\|_{\infty, K} \int_K |f(x)| dx$$
  Or, la convergence uniforme implique $\|\phi_j\|_{\infty, K} \to 0$, et comme $f \in L^1_{loc}$, l'intégrale sur $K$ est finie. Donc $\langle T_f, \phi_j \rangle \to 0$.
Toute fonction (assez régulière) peut donc être identifiée à une distribution.

\begin{quote}\textbf{Définition 5 (Convergence au sens des distributions) :} Une suite de distributions \end{quote}$(T_j)_{j \in \mathbb{N}}$ de $\mathcal{D}'(\mathbb{R}^n)$ converge vers une distribution $T \in \mathcal{D}'(\mathbb{R}^n)$ si, pour toute fonction test $\phi \in \mathcal{D}(\mathbb{R}^n)$ :
\begin{quote}\end{quote}$$\lim_{j \to +\infty} \langle T_j, \phi \rangle = \langle T, \phi \rangle$$
\begin{quote}C'est une convergence simple (dite convergence faible-*).\end{quote}

\textbf{Exemple 3 : Convergence vers le Dirac}
Posons la fonction porte (illustrée en Figure 1) : $f_j(x) = j$ si $x \in [-1/(2j), 1/(2j)]$ et $f_j(x) = 0$ sinon.
Calculons l'action de $T_{f_j}$ sur une fonction test $\phi \in \mathcal{D}(\mathbb{R})$ :
$$\langle T_{f_j}, \phi \rangle = \int_{\mathbb{R}} f_j(x)\phi(x) dx = j \int_{-1/(2j)}^{1/(2j)} \phi(x) dx$$
Par le théorème de la moyenne pour les intégrales, il existe $c_j \in [-1/(2j), 1/(2j)]$ tel que l'intégrale vaut la largeur du domaine $(1/j)$ multipliée par la valeur au point $c_j$.
$$\langle T_{f_j}, \phi \rangle = j \times \left( \frac{1}{j} \phi(c_j) \right) = \phi(c_j)$$
Lorsque $j \to +\infty$, l'intervalle se resserre vers $0$, donc $c_j \to 0$. Puisque $\phi$ est continue en $0$, $\phi(c_j) \to \phi(0) = \langle \delta_0, \phi \rangle$.
Ainsi, au sens des distributions, $\lim_{j \to +\infty} T_{f_j} = \delta_0$.

\textbf{Exemple 4 : La Valeur Principale de Cauchy}
La fonction $x \mapsto 1/x$ n'est pas localement intégrable autour de $0$, elle ne définit donc pas une distribution régulière. On peut définir une distribution singulière appelée "Valeur principale", notée $\text{vp}(1/x)$, par la limite d'intégrales évitant la singularité :
$$\langle \text{vp}\left(\frac{1}{x}\right), \phi \rangle = \lim_{\epsilon \to 0^+} \int_{|x| > \epsilon} \frac{\phi(x)}{x} dx$$

\section{Démonstrations}

\subsection{Proposition : Le Dirac n'est pas une distribution régulière}
Nous allons démontrer de manière exhaustive, par l'absurde, qu'il n'existe aucune fonction $f \in L^1_{loc}(\mathbb{R})$ qui puisse représenter la distribution de Dirac $\delta_0$.

\textbf{Démonstration :}
Supposons par l'absurde qu'il existe une fonction $f \in L^1_{loc}(\mathbb{R})$ telle que $T_f = \delta_0$, ce qui signifie que pour toute fonction test $\phi \in \mathcal{D}(\mathbb{R})$, nous avons :
$$\int_{-\infty}^{+\infty} f(x) \phi(x) dx = \phi(0)$$

Étape 1 : Construction d'une suite de fonctions tests appropriée.
Considérons une fonction test "bosse" standard, notée $\theta \in \mathcal{D}(\mathbb{R})$, possédant les propriétés suivantes :
- $0 \le \theta(x) \le 1$ pour tout $x \in \mathbb{R}$.
- $\theta(0) = 1$.
- $\text{supp}(\theta) \subset [-1, 1]$.
À partir de $\theta$, nous construisons une suite de fonctions tests concentrées autour de $0$ :
$$\phi_j(x) = \theta(jx)$$
Vérifions les propriétés de $\phi_j$ :
- $\phi_j \in \mathcal{D}(\mathbb{R})$ car la composition par une dilatation conserve le caractère $\mathcal{C}^\infty$ et le support compact.
- $\phi_j(0) = \theta(0) = 1$ pour tout $j$.
- Le support de $\phi_j$ est l'ensemble des points $x$ tels que $jx \in [-1, 1]$, soit $x \in [-1/j, 1/j]$.

Étape 2 : Évaluation des deux membres de notre égalité hypothétique pour $\phi_j$.
D'une part, puisque $T_f = \delta_0$, on a par définition de l'action de $\delta_0$ sur $\phi_j$ :
$$\langle \delta_0, \phi_j \rangle = \phi_j(0) = 1$$
Et ceci est strictement égal à 1 pour tout entier $j \ge 1$.

D'autre part, en utilisant la forme intégrale avec notre fonction $f \in L^1_{loc}(\mathbb{R})$ :
$$\langle T_f, \phi_j \rangle = \int_{-\infty}^{+\infty} f(x) \phi_j(x) dx = \int_{-1/j}^{1/j} f(x) \phi_j(x) dx$$
(Puisque $\phi_j(x) = 0$ en dehors de $[-1/j, 1/j]$).

Prenons la valeur absolue de cette intégrale et majorons-la :
$$\left| \int_{-1/j}^{1/j} f(x) \phi_j(x) dx \right| \le \int_{-1/j}^{1/j} |f(x)| |\phi_j(x)| dx$$
Étant donné que $0 \le \phi_j(x) \le 1$ pour tout $x$, on peut majorer par :
$$\left| \langle T_f, \phi_j \rangle \right| \le \int_{-1/j}^{1/j} |f(x)| dx$$

Étape 3 : Passage à la limite quand $j \to +\infty$.
La fonction $f$ étant localement intégrable, la fonction $|f|$ est une fonction positive dont l'intégrale sur le compact $[-1, 1]$ est finie.
Par le théorème de continuité absolue de l'intégrale de Lebesgue, si on intègre une telle fonction sur un domaine dont la mesure de Lebesgue tend vers $0$, la valeur de l'intégrale tend vers $0$.
Ici, l'intervalle d'intégration est $[-1/j, 1/j]$, dont la mesure (la longueur) est $2/j$, qui tend rigoureusement vers $0$ lorsque $j \to +\infty$.
Par conséquent :
$$\lim_{j \to +\infty} \int_{-1/j}^{1/j} |f(x)| dx = 0$$
Ce qui implique :
$$\lim_{j \to +\infty} \langle T_f, \phi_j \rangle = 0$$

Étape 4 : Conclusion.
Nous avons obtenu, par la représentation intégrale, une limite valant $0$. Mais par définition de l'égalité supposée avec $\delta_0$, cette même quantité devait valoir exactement $1$ pour tout $j$. Nous obtenons $0 = 1$, ce qui est une contradiction manifeste. L'hypothèse de départ est fausse : la masse de Dirac ne peut pas être représentée par une fonction localement intégrable. C'est un objet mathématique intrinsèquement différent, une "distribution singulière". $\blacksquare$

\section{Applications en Intelligence Artificielle}

En intelligence artificielle, et plus particulièrement en apprentissage statistique (Machine Learning), la théorie des distributions intervient de manière structurelle, bien que souvent implicite dans l'écriture des équations au quotidien.

\textbf{Modélisation des densités de probabilité empiriques :}
Lorsqu'un réseau de neurones est entraîné par apprentissage supervisé, le jeu de données d'entraînement est constitué d'un ensemble fini de couples $(x_i, y_i)$. La distribution sous-jacente des données génératrices n'est pas connue. L'algorithme d'optimisation (par exemple la descente de gradient stochastique) tente de minimiser l'espérance de la perte (le risque réel) définie formellement par rapport à une mesure de probabilité $P$ :
$$R(w) = \int \mathcal{L}(f_w(x), y) dP(x,y)$$
Mais $P$ étant inconnue, on approxime cette intégrale par l'espérance sur une mesure empirique $\hat{P}$, qui est une combinaison convexe de distributions de Dirac centrées sur les points d'entraînement :
$$\hat{P} = \frac{1}{N} \sum_{i=1}^N \delta_{(x_i, y_i)}$$
La mesure de probabilité empirique est donc une distribution au sens de Schwartz. Son action sur la fonction de perte $\mathcal{L}$ (agissant comme la "fonction test" dans ce contexte d'intégration de Lebesgue généralisée) redonne bien la moyenne arithmétique de la perte sur l'échantillon, le risque empirique que l'on minimise effectivement.

\textbf{Calcul des gradients aux singularités :}
Les fonctions d'activation modernes, telles que le ReLU ($f(x) = \max(0, x)$), présentent des points de non-dérivabilité classique (en $0$). Les algorithmes de rétropropagation calculent implicitement des dérivées au sens des distributions. La dérivée d'une fonction échelon (de Heaviside) est rigoureusement définie dans le cadre des distributions comme étant la distribution de Dirac. L'élargissement de l'espace des fonctions aux distributions est la fondation théorique autorisant le succès massif des réseaux profonds non lisses optimisés par sous-gradients.

\newpage
\part{Exercices d'Application et de Concours}
\subsection*{Exercice 1 : Action d'une distribution sur une fonction test de base \quad $\bigstar\star\star\star\star$}

Soit la distribution de Dirac $\delta_1$ et la fonction test $\phi(x) = \exp(-x^2)$.
1. Vérifier que $\phi$ n'est pas strictement dans $\mathcal{D}(\mathbb{R})$. Quelle restriction appliquer pour que $\langle \delta_1, \phi \rangle$ ait un sens rigoureux ?
2. Calculer l'action de $\delta_1$ sur une fonction test modifiée $\tilde{\phi}$ qui coïncide avec $\phi$ sur $[-2, 2]$ et est nulle en dehors de $[-3, 3]$.

\textbf{Correction Détaillée :}
1. La fonction $\phi(x) = \exp(-x^2)$ est infiniment dérivable sur $\mathbb{R}$ car la fonction exponentielle et les polynômes le sont (composition de fonctions $\mathcal{C}^\infty$). Cependant, $\phi(x) > 0$ pour tout $x \in \mathbb{R}$. Son support n'est donc pas compact, ce n'est pas une fonction test de $\mathcal{D}(\mathbb{R})$. $\delta_1$ étant une distribution, elle agit par définition sur $\mathcal{D}(\mathbb{R})$.
2. On définit une nouvelle fonction test $\tilde{\phi} \in \mathcal{D}(\mathbb{R})$ telle que $\tilde{\phi}(x) = \exp(-x^2)$ pour $x \in [-2, 2]$. L'action de la distribution de Dirac en $1$ sur une fonction test est l'évaluation de cette fonction au point $1$.
Puisque $1 \in [-2, 2]$, on a $\tilde{\phi}(1) = \exp(-1^2) = e^{-1}$.
Par conséquent, $\langle \delta_1, \tilde{\phi} \rangle = e^{-1}$. $\blacksquare$


\subsection*{Exercice 2 : La fonction de Heaviside \quad $\bigstar\star\star\star\star$}

Soit la fonction de Heaviside $H(x) = 1$ si $x \ge 0$ et $H(x) = 0$ si $x < 0$.
1. Justifier que $H$ définit une distribution régulière $T_H$.
2. Évaluer $\langle T_H, \phi \rangle$ pour une fonction test générale $\phi \in \mathcal{D}(\mathbb{R})$.

\textbf{Correction Détaillée :}
1. Pour qu'une fonction $f$ définisse une distribution régulière, il faut et il suffit qu'elle soit localement intégrable, c'est-à-dire $f \in L^1_{loc}(\mathbb{R})$.
Pour tout compact $K = [a, b]$ de $\mathbb{R}$, on a :
$\int_K |H(x)| dx = \int_a^b |H(x)| dx$.
Puisque $0 \le H(x) \le 1$, l'intégrale est majorée par $b-a$, qui est fini. Donc $H \in L^1_{loc}(\mathbb{R})$ et $H$ définit bien une distribution régulière $T_H$.
2. L'action de $T_H$ sur une fonction test $\phi \in \mathcal{D}(\mathbb{R})$ est donnée par :
$\langle T_H, \phi \rangle = \int_{-\infty}^{+\infty} H(x) \phi(x) dx$
Comme $H(x) = 0$ sur $]-\infty, 0[$ et $H(x) = 1$ sur $[0, +\infty[$, on obtient :
$\langle T_H, \phi \rangle = \int_0^{+\infty} 1 \cdot \phi(x) dx = \int_0^{+\infty} \phi(x) dx$.
Puisque $\phi$ est à support compact, l'intégrale converge (elle se réduit en fait à une intégrale sur un segment fermé et borné). $\blacksquare$


\subsection*{Exercice 3 : Linéarité des distributions singulières \quad $\bigstar\bigstar\star\star\star$}

On définit $T = 3\delta_{-2} - 5\delta_{4}$.
1. Montrer que $T$ est une distribution.
2. Calculer $\langle T, \phi \rangle$ où $\phi(x) = x^2$ restreinte de manière adéquate à un compact.

\textbf{Correction Détaillée :}
1. L'espace $\mathcal{D}'(\mathbb{R})$ est un espace vectoriel. Puisque $\delta_{-2}$ et $\delta_{4}$ sont des distributions (le Dirac est une forme linéaire continue), toute combinaison linéaire scalaire l'est également. Ainsi, la somme $T = 3\delta_{-2} - 5\delta_{4}$ est une forme linéaire continue sur $\mathcal{D}(\mathbb{R})$, c'est donc une distribution.
2. Soit $\tilde{\phi}$ une fonction de $\mathcal{D}(\mathbb{R})$ telle que $\tilde{\phi}(x) = x^2$ sur un intervalle compact contenant à la fois $-2$ et $4$, par exemple $[-5, 5]$.
On calcule l'action de $T$ sur $\tilde{\phi}$ par linéarité :
$\langle T, \tilde{\phi} \rangle = \langle 3\delta_{-2} - 5\delta_{4}, \tilde{\phi} \rangle$
$\langle T, \tilde{\phi} \rangle = 3 \langle \delta_{-2}, \tilde{\phi} \rangle - 5 \langle \delta_{4}, \tilde{\phi} \rangle$
En utilisant la définition du Dirac, l'action est l'évaluation au point :
$\langle T, \tilde{\phi} \rangle = 3 \tilde{\phi}(-2) - 5 \tilde{\phi}(4)$
Puisque $\tilde{\phi}(x) = x^2$ en ces points :
$\langle T, \tilde{\phi} \rangle = 3(-2)^2 - 5(4)^2 = 3(4) - 5(16) = 12 - 80 = -68$. $\blacksquare$


\subsection*{Exercice 4 : Translation d'une distribution régulière \quad $\bigstar\bigstar\star\star\star$}

Soit $f \in L^1_{loc}(\mathbb{R})$ et $T_f$ sa distribution régulière associée. On définit la fonction translatée $\tau_a f(x) = f(x - a)$ pour $a \in \mathbb{R}$.
1. Calculer l'action de $T_{\tau_a f}$ sur une fonction test $\phi$.
2. En déduire une définition générale pour la translatée d'une distribution quelconque $\tau_a T$, et l'appliquer à $\delta_0$.

\textbf{Correction Détaillée :}
1. L'action de la distribution régulière associée à $f(x-a)$ est :
$\langle T_{\tau_a f}, \phi \rangle = \int_{-\infty}^{+\infty} f(x - a) \phi(x) dx$
Effectuons le changement de variable affine (donc bijectif et $\mathcal{C}^\infty$) $y = x - a$. On a $dx = dy$ :
$\langle T_{\tau_a f}, \phi \rangle = \int_{-\infty}^{+\infty} f(y) \phi(y + a) dy$
On peut réécrire cette intégrale en utilisant l'action de $T_f$ :
$\langle T_{\tau_a f}, \phi \rangle = \langle T_f, \phi(\cdot + a) \rangle$.
Ou en posant $\tau_{-a}\phi(y) = \phi(y+a)$ :
$\langle T_{\tau_a f}, \phi \rangle = \langle T_f, \tau_{-a}\phi \rangle$.
2. Cette formule n'utilise plus explicitement d'intégrale de Lebesgue, elle dépend uniquement de l'action de la distribution sur la fonction test translatée (qui est toujours dans $\mathcal{D}(\mathbb{R})$). On peut donc généraliser et définir la translatée d'une distribution $T$ par :
$\langle \tau_a T, \phi \rangle = \langle T, \tau_{-a}\phi \rangle$
Appliquons cela à $\delta_0$ :
$\langle \tau_a \delta_0, \phi \rangle = \langle \delta_0, \phi(\cdot + a) \rangle$
L'action de $\delta_0$ sur une fonction est l'évaluation de cette fonction en $0$ :
$\langle \tau_a \delta_0, \phi \rangle = \phi(0 + a) = \phi(a)$
Or $\phi(a) = \langle \delta_a, \phi \rangle$.
On conclut que la translatée de $\delta_0$ par un vecteur $a$ est $\delta_a$ : $\tau_a \delta_0 = \delta_a$. $\blacksquare$


\subsection*{Exercice 5 : Convergence de suite de fonctions vers un Dirac \quad $\bigstar\bigstar\bigstar\star\star$}

Soit la suite de fonctions $f_n(x) = \frac{n}{\pi (1 + n^2 x^2)}$.
Montrer que $f_n \to \delta_0$ au sens des distributions quand $n \to +\infty$.

\textbf{Correction Détaillée :}
Soit $\phi \in \mathcal{D}(\mathbb{R})$. Nous devons montrer que $\lim_{n \to +\infty} \int_{-\infty}^{+\infty} f_n(x) \phi(x) dx = \phi(0)$.
Séparons l'intégrale en deux parties (on injecte $\phi(0)$) :
$\int_{-\infty}^{+\infty} f_n(x) \phi(x) dx = \int_{-\infty}^{+\infty} \frac{n}{\pi (1 + n^2 x^2)} (\phi(x) - \phi(0) + \phi(0)) dx$
$ = \phi(0) \int_{-\infty}^{+\infty} \frac{n}{\pi (1 + n^2 x^2)} dx + \int_{-\infty}^{+\infty} \frac{n}{\pi (1 + n^2 x^2)} (\phi(x) - \phi(0)) dx$

Calculons la première intégrale (on effectue $y = nx$, $dy = n dx$) :
$\int_{-\infty}^{+\infty} \frac{n}{\pi (1 + n^2 x^2)} dx = \frac{1}{\pi} \int_{-\infty}^{+\infty} \frac{1}{1 + y^2} dy = \frac{1}{\pi} [\arctan(y)]_{-\infty}^{+\infty} = \frac{1}{\pi} \left( \frac{\pi}{2} - \left(-\frac{\pi}{2}\right) \right) = 1$.
La première partie vaut donc exactement $\phi(0)$.

Il reste à montrer que le second terme tend vers 0. Soit $I_n = \int_{-\infty}^{+\infty} \frac{n}{\pi (1 + n^2 x^2)} (\phi(x) - \phi(0)) dx$.
Par le théorème des accroissements finis, puisque $\phi \in \mathcal{C}^\infty$ et est à support compact, sa dérivée $\phi'$ est bornée sur $\mathbb{R}$. Soit $M = \sup_{\mathbb{R}} |\phi'(x)|$.
On a $|\phi(x) - \phi(0)| \le M |x|$ pour tout $x$.
On scinde l'intégrale en $|x| \le A/\sqrt{n}$ et $|x| > A/\sqrt{n}$ pour un $A>0$ bien choisi.
Alternativement (plus rapide), on effectue le changement de variable $y = nx$ :
$I_n = \frac{1}{\pi} \int_{-\infty}^{+\infty} \frac{1}{1 + y^2} \left( \phi\left(\frac{y}{n}\right) - \phi(0) \right) dy$
Pour chaque $y$ fixé, la fonction intégrande tend vers $0$ car $\phi$ est continue en $0$.
De plus, la fonction intégrande est majorée en valeur absolue par $\frac{2 \|\phi\|_\infty}{1+y^2}$, qui est une fonction dans $L^1(\mathbb{R})$ indépendante de $n$.
On peut donc appliquer le Théorème de Convergence Dominée de Lebesgue : l'intégrale de la limite est la limite de l'intégrale, ce qui donne $0$.
Ainsi, $\lim_{n \to +\infty} \langle T_{f_n}, \phi \rangle = \phi(0) + 0 = \langle \delta_0, \phi \rangle$.
La suite converge bien vers $\delta_0$. $\blacksquare$


\subsection*{Exercice 6 : Multiplication par une fonction de classe C-infini \quad $\bigstar\bigstar\bigstar\star\star$}

Soit $T \in \mathcal{D}'(\mathbb{R})$ et $\alpha \in \mathcal{C}^\infty(\mathbb{R})$.
1. Comment définir l'action de la distribution $\alpha T$ sur une fonction test $\phi$ ?
2. Montrer que l'équation $x T = 0$ admet pour solution générale $T = c \delta_0$ (où $c$ est une constante complexe).

\textbf{Correction Détaillée :}
1. Si $T$ était une distribution régulière associée à $f \in L^1_{loc}$, l'action serait :
$\langle \alpha T_f, \phi \rangle = \int (\alpha(x) f(x)) \phi(x) dx = \int f(x) (\alpha(x) \phi(x)) dx$.
Pour que cette forme intégrale s'étende aux distributions, il faut que $\alpha \phi$ soit une fonction test. Comme $\phi \in \mathcal{D}(\mathbb{R})$ (elle est $\mathcal{C}^\infty$ à support compact) et $\alpha \in \mathcal{C}^\infty(\mathbb{R})$, le produit $\alpha \phi$ reste $\mathcal{C}^\infty$ et le support de $\alpha \phi$ est inclus dans le support de $\phi$, donc compact.
On définit formellement : $\langle \alpha T, \phi \rangle = \langle T, \alpha \phi \rangle$.
2. L'équation est $x T = 0$, ce qui signifie que pour tout $\phi \in \mathcal{D}(\mathbb{R})$, $\langle x T, \phi \rangle = 0$, soit $\langle T, x\phi \rangle = 0$.
Soit une fonction test $\phi$. On peut écrire $\phi(x) = \phi(0)\theta(x) + x \psi(x)$ où $\theta \in \mathcal{D}(\mathbb{R})$ est une fonction test valant $1$ au voisinage de $0$, et $\psi \in \mathcal{D}(\mathbb{R})$.
En effet, $\psi(x) = \frac{\phi(x) - \phi(0)\theta(x)}{x}$. Par le théorème de Taylor, cette fonction est bien prolongeable par continuité (et régularité $\mathcal{C}^\infty$) en $0$.
Appliquons $T$ :
$\langle T, \phi \rangle = \langle T, \phi(0)\theta \rangle + \langle T, x\psi \rangle$
Par hypothèse, $\langle T, x\psi \rangle = \langle xT, \psi \rangle = 0$.
Ainsi, $\langle T, \phi \rangle = \phi(0) \langle T, \theta \rangle$.
La quantité $\langle T, \theta \rangle$ est une constante fixe que nous noterons $c$.
D'où $\langle T, \phi \rangle = c \phi(0) = \langle c \delta_0, \phi \rangle$.
La solution est bien $T = c \delta_0$. $\blacksquare$


\subsection*{Exercice 7 : La distribution Valeur Principale (vp) \quad $\bigstar\bigstar\bigstar\bigstar\star$}

On définit pour $\phi \in \mathcal{D}(\mathbb{R})$ la forme linéaire $\text{vp}(1/x)$ par :
$$\langle \text{vp}\left(\frac{1}{x}\right), \phi \rangle = \lim_{\epsilon \to 0^+} \int_{|x|>\epsilon} \frac{\phi(x)}{x} dx$$
Montrer que cette limite existe bien pour toute fonction test, démontrant ainsi qu'il s'agit d'une distribution.

\textbf{Correction Détaillée :}
Pour contourner la divergence au voisinage de zéro due à $1/x$, on exploite la symétrie.
L'intégrale se sépare sur deux domaines :
$\int_{|x|>\epsilon} \frac{\phi(x)}{x} dx = \int_{-\infty}^{-\epsilon} \frac{\phi(x)}{x} dx + \int_{\epsilon}^{+\infty} \frac{\phi(x)}{x} dx$
Dans la première intégrale, on effectue le changement de variable $x \to -x$ :
$\int_{-\infty}^{-\epsilon} \frac{\phi(x)}{x} dx = \int_{+\infty}^{\epsilon} \frac{\phi(-x)}{-x} (-dx) = -\int_{\epsilon}^{+\infty} \frac{\phi(-x)}{x} dx$
En recombinant avec la seconde intégrale :
$\int_{|x|>\epsilon} \frac{\phi(x)}{x} dx = \int_{\epsilon}^{+\infty} \frac{\phi(x) - \phi(-x)}{x} dx$
Par le théorème des accroissements finis appliqué entre $-x$ et $x$, $\phi(x) - \phi(-x) = 2x \phi'(c_x)$ avec $c_x \in [-x, x]$.
Puisque $\phi$ est $\mathcal{C}^\infty$, sa dérivée est bornée, disons par $M$.
Donc $\left| \frac{\phi(x) - \phi(-x)}{x} \right| \le \frac{2|x|M}{|x|} = 2M$.
La fonction intégrande $\frac{\phi(x) - \phi(-x)}{x}$ est bornée au voisinage de $0$.
De plus, son intégration se fait sur un ensemble compact puisque $\phi$ est à support compact. L'intégrale $\int_{0}^{+\infty} \frac{\phi(x) - \phi(-x)}{x} dx$ est donc parfaitement convergente (absolument convergente).
La limite quand $\epsilon \to 0^+$ existe donc bel et bien. $\text{vp}(1/x)$ est une distribution bien définie. $\blacksquare$


\subsection*{Exercice 8 : Distribution issue du logarithme \quad $\bigstar\bigstar\bigstar\bigstar\star$}

On note $f(x) = \ln|x|$.
1. Montrer que $f \in L^1_{loc}(\mathbb{R})$. En déduire que $T_f$ est une distribution.
2. (Préparation au jalon suivant) Si l'on calculait formellement $\langle T_f', \phi \rangle = - \langle T_f, \phi' \rangle$, quelle distribution singulière reconnaîtrait-on ?

\textbf{Correction Détaillée :}
1. La fonction logarithme est mesurable. Son seul point de singularité est $0$. Étudions l'intégrabilité locale en $0$.
Il s'agit de voir si l'intégrale $\int_{-\epsilon}^{\epsilon} |\ln|x|| dx$ converge pour $\epsilon > 0$.
Par symétrie, l'intégrale vaut $2 \int_0^{\epsilon} (-\ln(x)) dx$.
Une primitive de $\ln(x)$ est $x \ln(x) - x$.
En évaluant de $\delta > 0$ à $\epsilon$ puis en passant à la limite $\delta \to 0^+$ :
$\lim_{\delta \to 0^+} [\delta \ln(\delta) - \delta] = 0$. (Croissance comparée usuelle).
L'intégrale vaut donc $2(\epsilon - \epsilon\ln(\epsilon))$, ce qui est fini.
Ainsi, $f$ est localement intégrable et $T_f$ est une distribution bien définie (régulière).
2. Calculons $-\langle T_f, \phi' \rangle$ :
$-\langle T_f, \phi' \rangle = - \int_{-\infty}^{+\infty} \ln|x| \phi'(x) dx = - \lim_{\epsilon \to 0^+} \left( \int_{-\infty}^{-\epsilon} \ln(-x) \phi'(x) dx + \int_{\epsilon}^{+\infty} \ln(x) \phi'(x) dx \right)$
Intégrons par parties sur les domaines ne contenant pas 0. Le terme de bord à l'infini s'annule car $\phi$ est à support compact.
$- \left[ \ln(-x)\phi(x) \right]_{-\infty}^{-\epsilon} + \int_{-\infty}^{-\epsilon} \frac{\phi(x)}{x} dx - \left[ \ln(x)\phi(x) \right]_{\epsilon}^{+\infty} + \int_{\epsilon}^{+\infty} \frac{\phi(x)}{x} dx$
$- \ln(\epsilon)\phi(-\epsilon) + \ln(\epsilon)\phi(\epsilon) + \int_{|x| > \epsilon} \frac{\phi(x)}{x} dx$
Le terme de bord s'écrit $\ln(\epsilon)(\phi(\epsilon) - \phi(-\epsilon))$. Or $\phi(\epsilon) - \phi(-\epsilon) \sim 2\epsilon \phi'(0)$ quand $\epsilon \to 0$.
Donc $\ln(\epsilon) \times 2\epsilon \phi'(0) \to 0$ par croissance comparée.
Il reste l'intégrale, qui, à la limite, n'est autre que la définition de la valeur principale.
On trouve ainsi formellement que la dérivée du logarithme au sens des distributions est $\text{vp}(1/x)$. $\blacksquare$


\subsection*{Exercice 9 : Égalité de distributions : vp et partie réelle \quad $\bigstar\bigstar\bigstar\bigstar\bigstar$}

Soit $\epsilon > 0$ et $f_\epsilon(x) = \frac{1}{x + i\epsilon}$.
Montrer que, au sens des distributions (c'est-à-dire la limite de l'action de $f_\epsilon$ sur une fonction test quand $\epsilon \to 0^+$) :
$\lim_{\epsilon \to 0^+} \frac{1}{x + i\epsilon} = \text{vp}\left(\frac{1}{x}\right) - i \pi \delta_0$
(C'est l'identité de Sokhotski-Plemelj).

\textbf{Correction Détaillée :}
Séparons la partie réelle et la partie imaginaire de $f_\epsilon$ :
$\frac{1}{x + i\epsilon} = \frac{x - i\epsilon}{x^2 + \epsilon^2} = \frac{x}{x^2 + \epsilon^2} - i \frac{\epsilon}{x^2 + \epsilon^2}$
Soit $\phi \in \mathcal{D}(\mathbb{R})$.
Partie imaginaire : $-\int_{-\infty}^{+\infty} \frac{\epsilon}{x^2 + \epsilon^2} \phi(x) dx$.
On reconnaît une approximation de l'identité proportionnelle au noyau de Poisson (cf. Ex 5).
Posons $y = x/\epsilon$. L'intégrale devient $-\int_{-\infty}^{+\infty} \frac{\epsilon}{\epsilon^2 y^2 + \epsilon^2} \phi(\epsilon y) \epsilon dy = - \int_{-\infty}^{+\infty} \frac{1}{y^2 + 1} \phi(\epsilon y) dy$.
Par convergence dominée (la fonction est dominée par $C/(1+y^2)$ et $\phi(\epsilon y) \to \phi(0)$), l'intégrale tend vers $-\int_{-\infty}^{+\infty} \frac{1}{y^2+1} \phi(0) dy = -\pi \phi(0) = \langle -i\pi \delta_0, \phi \rangle$.
Partie réelle : $\int_{-\infty}^{+\infty} \frac{x}{x^2 + \epsilon^2} \phi(x) dx$.
Le terme $\frac{x}{x^2 + \epsilon^2}$ est impair. On peut donc injecter $\phi(0)$ en remarquant que son intégrale avec une fonction impaire sur un domaine symétrique est nulle :
$\int_{-\infty}^{+\infty} \frac{x}{x^2 + \epsilon^2} (\phi(x) - \phi(0)) dx$.
La fonction $\psi(x) = (\phi(x)-\phi(0))/x$ est lisse et à support compact (prolongeable en 0).
L'intégrale s'écrit $\int_{-\infty}^{+\infty} \frac{x^2}{x^2+\epsilon^2} \psi(x) dx$.
Quand $\epsilon \to 0^+$, l'intégrande converge vers $\psi(x)$ presque partout et est dominé.
La limite est donc $\int_{-\infty}^{+\infty} \psi(x) dx = \int_{-\infty}^{+\infty} \frac{\phi(x)-\phi(0)}{x} dx$, qui est l'expression équivalente de $\text{vp}(1/x)$ démontrée dans l'exercice 7.
La combinaison des deux parties donne la relation de Sokhotski-Plemelj. $\blacksquare$


\subsection*{Exercice 10 : L'équation différentielle $x T' = 0$ \quad $\bigstar\bigstar\bigstar\bigstar\bigstar$}

(Exercice d'ouverture vers le Jalon 83).
Déterminer toutes les distributions $T \in \mathcal{D}'(\mathbb{R})$ telles que $x T' = 0$, sachant que par définition l'action de $T'$ est $\langle T', \phi \rangle = - \langle T, \phi' \rangle$.

\textbf{Correction Détaillée :}
D'après l'exercice 6, si $S \in \mathcal{D}'$ vérifie $xS = 0$, alors $S = c \delta_0$ pour une constante $c$.
Ici on a $x T' = 0$, donc $T' = c_1 \delta_0$.
Il faut alors chercher les primitives de la distribution de Dirac.
Par définition de la dérivée, on veut : $\langle T, -\phi' \rangle = c_1 \phi(0)$ pour toute fonction test.
Soit $\psi$ une fonction test qui s'écrit sous la forme d'une dérivée : $\psi = \phi'$ pour une certaine $\phi \in \mathcal{D}(\mathbb{R})$.
Une fonction test $\psi$ est la dérivée d'une fonction test $\phi$ (i.e. $\phi(x) = \int_{-\infty}^x \psi(t)dt \in \mathcal{D}(\mathbb{R})$) si et seulement si l'intégrale totale de $\psi$ est nulle : $\int_{-\infty}^{+\infty} \psi(t) dt = 0$.
On sait d'autre part (Exercice 1) que la dérivée de la fonction de Heaviside est le Dirac. Prouvons-le rigoureusement : $\langle H', \phi \rangle = - \langle H, \phi' \rangle = - \int_0^{+\infty} \phi'(x) dx = - (0 - \phi(0)) = \phi(0) = \langle \delta_0, \phi \rangle$.
Donc $H$ est UNE solution.
Par linéarité de la dérivation, la solution générale sera de la forme $c_1 H + c_2$ (une constante). Une distribution constante est la distribution régulière associée à une fonction constante $f(x) = c_2$.
Vérifions : $T = c_1 T_H + c_2 T_1$.
$\langle (c_1 T_H + c_2 T_1)', \phi \rangle = -c_1 \langle T_H, \phi' \rangle - c_2 \int 1 \cdot \phi'(x) dx$.
L'intégrale de $\phi'$ est nulle car $\phi$ est à support compact.
On a bien $(c_1 T_H + c_2 T_1)' = c_1 \delta_0$.
Puis $x(c_1 \delta_0) = 0$.
L'ensemble des distributions solutions est donc l'ensemble des formes $T = c_1 T_H + c_2 T_1$, c'est-à-dire des "marches d'escalier" dont le saut est en 0. $\blacksquare$



\newpage
\part{Travaux Pratiques et Simulations Algorithmiques}
\subsection*{TP 1 : Simulation de la convergence vers la distribution de Dirac}

L'objectif de ce TP est de simuler numeriquement l'action d'une suite de fonctions (portes) convergeant vers la distribution de Dirac sur une fonction test choisie analytiquement, pour verifier que l'integrale numerique converge bien vers la valeur de la fonction en zero.

\begin{lstlisting}[language=Python]
import math

def fonction_test_cloche(x):
    # Fonction lisse a support compact [-1, 1]
    if abs(x) >= 1.0:
        return 0.0
    return math.exp(-1.0 / (1.0 - x**2))

def fonction_porte(x, n):
    # Porte d'aire 1 et de largeur 1/n centree en 0
    if abs(x) <= 1.0 / (2.0 * n):
        return n
    return 0.0

def integration_trapezes(f, a, b, N=10000):
    h = (b - a) / N
    s = 0.5 * (f(a) + f(b))
    for i in range(1, N):
        s += f(a + i * h)
    return s * h

def action_sur_test(n):
    # L'integrale se limite au support de la porte [-1/(2n), 1/(2n)]
    a = -1.0 / (2.0 * n)
    b = 1.0 / (2.0 * n)
    # L'integrande est porte(x)*test(x)
    integrande = lambda x: fonction_porte(x, n) * fonction_test_cloche(x)
    return integration_trapezes(integrande, a, b)

valeur_theorique = fonction_test_cloche(0.0)
print(f"Valeur theorique (Dirac en 0) : {valeur_theorique}")

for n in [1, 10, 100, 1000]:
    valeur_approx = action_sur_test(n)
    erreur = abs(valeur_approx - valeur_theorique)
    print(f"n={n:4d} | Approx: {valeur_approx:.6f} | Erreur: {erreur:.6e}")
    # Verification mathematique de la convergence
    assert erreur < 2.0 / n, "La convergence est trop lente ou incorrecte"

print("Convergence validee : la suite de portes agit comme un Dirac.")
\end{lstlisting}


\subsection*{TP 2 : Action de la translatée d'une distribution}

Nous allons implementer et verifier numeriquement la formule de translation d'une distribution. La translatee du Dirac en 0 par un vecteur 'a' doit etre le Dirac en 'a'.

\begin{lstlisting}[language=Python]
import math

def fonction_test_sinus(x):
    # Support [-math.pi, math.pi]
    if abs(x) >= math.pi:
        return 0.0
    return math.sin(x) * math.exp(-1.0 / (math.pi**2 - x**2))

def fonction_gaussienne_dirac(x, epsilon):
    # Approximation lisse du Dirac d'aire 1
    return (1.0 / (math.sqrt(2 * math.pi) * epsilon)) * math.exp(-0.5 * (x / epsilon)**2)

def integration_simpson(f, a, b, N=1000):
    if N % 2 != 0:
        N += 1
    h = (b - a) / N
    s = f(a) + f(b)
    for i in range(1, N, 2):
        s += 4 * f(a + i * h)
    for i in range(2, N-1, 2):
        s += 2 * f(a + i * h)
    return s * h / 3.0

# On decale la distribution de a = 1.5
a_trans = 1.5
epsilon = 0.01

# Methode 1 : Integrale de (f(x-a) * phi(x))
def integrande_1(x):
    return fonction_gaussienne_dirac(x - a_trans, epsilon) * fonction_test_sinus(x)

# Methode 2 : Integrale de (f(x) * phi(x+a))
def integrande_2(x):
    return fonction_gaussienne_dirac(x, epsilon) * fonction_test_sinus(x + a_trans)

# Le support utile pour l'integrale est centre autour du pic de la gaussienne
val1 = integration_simpson(integrande_1, a_trans - 5*epsilon, a_trans + 5*epsilon)
val2 = integration_simpson(integrande_2, -5*epsilon, +5*epsilon)
val_theorique = fonction_test_sinus(a_trans)

print(f"Action theorique (Dirac en {a_trans}) : {val_theorique:.6f}")
print(f"Action integrale directe : {val1:.6f}")
print(f"Action sur fonction translatee : {val2:.6f}")

assert abs(val1 - val2) < 1e-5, "Les deux formulations integrales doivent coincider"
assert abs(val1 - val_theorique) < 1e-3, "L'approximation gaussienne doit etre proche du Dirac"
\end{lstlisting}


\subsection*{TP 3 : Calcul numérique de la Valeur Principale (vp)}

La distribution vp(1/x) est definie par une limite en evitant la singularite. Ce TP calcule cette limite numeriquement sur une fonction test asymetrique.

\begin{lstlisting}[language=Python]
import math

def fonction_test_asymetrique(x):
    if abs(x) >= 2.0:
        return 0.0
    # Fonction lisse non nulle et non paire
    return (x**2 - 4)**2 * math.exp(x)

def integrande_vp(x):
    # La fonction vp : phi(x)/x. Pas de division par zero car evaluee hors de 0
    return fonction_test_asymetrique(x) / x

def integration_vp(epsilon, a=-2.0, b=2.0):
    # On integre sur [a, -epsilon] et [epsilon, b]
    N = 10000

    # Partie negative
    h1 = (-epsilon - a) / N
    s1 = 0.5 * (integrande_vp(a) + integrande_vp(-epsilon))
    for i in range(1, N):
        s1 += integrande_vp(a + i * h1)
    s1 *= h1

    # Partie positive
    h2 = (b - epsilon) / N
    s2 = 0.5 * (integrande_vp(epsilon) + integrande_vp(b))
    for i in range(1, N):
        s2 += integrande_vp(epsilon + i * h2)
    s2 *= h2

    return s1 + s2

epsilons = [0.1, 0.01, 0.001, 0.0001]
valeurs_vp = []

for eps in epsilons:
    val = integration_vp(eps)
    valeurs_vp.append(val)
    print(f"Epsilon = {eps:.4f} -> Action vp = {val:.5f}")

# Verification de la stabilite de la limite (Convergence de Cauchy)
diff = abs(valeurs_vp[-1] - valeurs_vp[-2])
assert diff < 1e-3, "La valeur principale ne converge pas numeriquement"
print("La limite de la valeur principale est stable.")
\end{lstlisting}


\subsection*{TP 4 : Preuve numérique de l'identité de Sokhotski-Plemelj}

Nous allons simuler l'identite 1/(x+iepsilon) = vp(1/x) - i pi Dirac(0) en evaluant l'action de ces distributions sur une fonction test.

\begin{lstlisting}[language=Python]
import math

def fonction_test(x):
    if abs(x) >= 1.0:
        return 0.0
    return math.exp(x) * (1 - x**2)**3

def integration_complexe(epsilon, N=10000):
    a, b = -1.0, 1.0
    h = (b - a) / N
    sum_real = 0.0
    sum_imag = 0.0

    for i in range(N + 1):
        x = a + i * h
        poids = 0.5 if (i == 0 or i == N) else 1.0
        phi_val = fonction_test(x)

        # Denominateur : x^2 + epsilon^2
        den = x**2 + epsilon**2
        real_part = (x / den) * phi_val
        imag_part = (-epsilon / den) * phi_val

        sum_real += poids * real_part
        sum_imag += poids * imag_part

    return sum_real * h, sum_imag * h

val_dirac = fonction_test(0.0)
theorique_imag = -math.pi * val_dirac

print(f"Partie imaginaire theorique (-pi * phi(0)) : {theorique_imag:.5f}")

for eps in [0.1, 0.01, 0.001]:
    r, c = integration_complexe(eps)
    print(f"Eps={eps:.3f} | Reelle (vp)={r:.5f} | Imaginaire={c:.5f}")
    if eps == 0.001:
        assert abs(c - theorique_imag) < 1e-2, "La partie imaginaire ne tend pas vers -pi * Dirac"

print("Identite verifiee avec succes au sens des distributions.")
\end{lstlisting}


\subsection*{TP 5 : Résolution matricielle de l'équation différentielle xT=0 discretisée}

Nous simulons le theoreme qui stipule que si xT = 0, alors T est proportionnel au Dirac. Nous discretisons l'espace et resolvons le systeme lineaire associe.

\begin{lstlisting}[language=Python]
# Pas de module externe complexe pour rester pur Python
# On construit une grille 1D de -1 a 1 avec N points
N = 11
x_grid = [ -1.0 + i * (2.0 / (N - 1)) for i in range(N) ]

# Une distribution T discretisee est un vecteur d'evaluations
# L'equation (x * T) = 0 devient x_i * T_i = 0 pour chaque i
# Pour un systeme lineaire simple, on regarde le noyau de l'operateur de multiplication.

T_solution = [0.0] * N

for i in range(N):
    x = x_grid[i]
    # Si x n'est pas nul, T_i doit etre nul pour que x_i * T_i = 0
    # A cause de la precision flottante, on regarde si |x| est tres petit
    if abs(x) > 1e-10:
        T_solution[i] = 0.0
    else:
        # En x=0, n'importe quelle valeur est solution (proportionnalite au Dirac)
        T_solution[i] = 1.0 # On fixe arbitrairement c=1

print("Grille X   :", [round(x, 2) for x in x_grid])
print("Solution T :", T_solution)

# Verification de l'equation
for i in range(N):
    produit = x_grid[i] * T_solution[i]
    assert abs(produit) < 1e-10, "L'equation xT=0 n'est pas respectee"

# On verifie que seul le point central (indice 5 pour N=11, correspondant a x=0) est non nul
index_zero = N // 2
assert T_solution[index_zero] == 1.0, "Le Dirac doit etre au centre"
for i in range(N):
    if i != index_zero:
        assert T_solution[i] == 0.0, "Les autres valeurs doivent etre nulles"

print("Le noyau de l'operateur de multiplication par X est bien engendre par le Dirac discret.")
\end{lstlisting}


\end{document}
